In this section, we combine the construction of chain complexes of flow categories with the construction of spectral sequences arising from chain complexes. Our goal is to construct a spectral sequence that converges to the realization of a framed flow category, where the pages and differentials are described in terms of geometric data extracted from the flow category.

\subsection{Toda brackets in flow categories}

We now define a notion analogous to the Toda bracket in the context of framed flow categories.

\begin{defn}\label{def:Toda_FFCat}
    Let $\fC, \sD$ be framed flow categories such that there is a index preserving bijection
    \begin{equation*}
        \mu: \ob(\fC_{p,p-r}) \simeq \ob(\sD_{[p, p-r]})
    \end{equation*}
    and for $u,v \in \ob(\fC_{p,p-r})$ we have
    \begin{equation*}
        M_{u,v}^{\fC} \simeq M_{\mu(u),\mu(v)}^{\sD}.
    \end{equation*}
    We denote this relation as $\fC_{[p,p-r]} \simeq \sD_{[p,p-r]}$.
    
    Using \cref{constr:FF_to_Ch}, we can view $\fC_{[p+1,p-r]}$ and $\sD_{[p,p-r-1]}$ as $(p+1-r)$-chain complexes, denoted by $f_{\fC},f_{\sD}$. Using \cref{cor:ch_to_morph}, we can view $f_{\fC},f_{\sD}$ as morphisms of $(p-r)$-chain complexes of the form:
    \begin{equation*}
        f_{\fC}: X_p \to K, \quad f_{\sD}:  K' \to Y_{p-r}.
    \end{equation*}
    The requirement $\fC_{[p,p-r]} \simeq \sD_{[p,p-r]}$ imply that $K = K'$. The composition of yields a morphism of chain complexes
    \begin{equation*}
        f_{\sD} \circ f_{\fC}:X_p \to Y_{p-r},
    \end{equation*}
    which by \cref{constr:FF_to_Ch} is a chain complex with two non trivial objects. Using \cref{prop:FF_are_Ch} we get a framed flow category with two objects and as we saw in \cref{ex:closed_flow_cat}, such category correspond to framed closed manifolds of dimension $r$ which we call the $(p,r)$ \textbf{Toda bracket} of $\fC$ and $\sD$, denoted by:
    \begin{equation*}
        \<{\fC_{[p+1,p-r]}, \sD_{[p,p-r-1]}} \in \Omega^{\fr}_r.
    \end{equation*}
    
    We will write $\langle \fC,\sD \rangle$ if it's clear what are $p,r$. For an element $y \in \sD_{p-r-1}$ we denote by 
    \begin{equation*}
        \<{\fC,\sD}_y \coloneqq \<{\fC_{[p+1,p-r]},\sD_{[p,p-r]} \sqcup \{y\}}
    \end{equation*}
    the Toda brackets of $\fC_{[p+1,p-r]}$ and the full subcategory generated by $\sD_{[p,p-r]}$ and $y$.
    
    We can extend the definition in the following way: if there is $n \in \bbZ$ such that $\fC_{k}$ is trivial for  $n+1 >  k > p-r$, and $\fC_{[p,p-r]} \simeq \sD_{[p,p-r]}$, we get a $n-p+r$ bordism type:
    \begin{equation*}
       \<{\fC_{[n+1,p-r]}, \sD_{[p,p-r]}}_y \coloneqq \<{\fC_{[n+1,p-r]},\sD_{[p,p-r]} \sqcup \{y\}} \in \Omega^{\fr}_{n-p+r}
    \end{equation*}
\end{defn}

Later we will show that this definition of Toda bracket coincidence with \cref{def:smallToda}.

\begin{rem}
    A priory, \cref{prop:FF_are_Ch} give a lift, from chain complexes to framed flow categories, so as a framed manifold $\langle \fC,\sD \rangle$ depends on the lift. Given two lifts $\fC,\fC'$ such that $K_{\fC} \simeq K_{\fC} \in \Ch(\Sp)$, we can use the facts that $K_{\fC}$ have two object, which are $\bbS$, to get that $K_{\fC}$ is determined by an element $\pi_r(\bbS) \simeq \Omega^{\fr}_r$, hence the Toda brackets are well define as an element of the cobordism ring.
\end{rem}

\subsection{The spectral sequence of a framed flow category}

We now state the main result, which describes the spectral sequence associated to a framed flow category.
We recall that 
\begin{thm}\label{thm:SSofFFCat}
Let $\fC$ be a framed flow category. Then there exists a spectral sequence converging to the realization
\begin{equation*}
    \rnp{E} \implies \pi_n \norm{\fC},
\end{equation*}
with:
\begin{equation*}
    \pg{E}{1}{n}{p} = \bigoplus_{\norm{y} = p} \Omega^{\fr}_{n - p}.
\end{equation*}
An element $X \in \pg{E}{1}{n}{p}$ survives to the $(r+1)$-th page if and only if there is an $(X,r,n,p)$ extension of $\fC$ (see \cref{def:Xrnp-extension}). Let $\fC'$ be such extension, then the differential is represented in $\pg{E}{1}{n - 1}{p - (r+1)}$ by the Toda bracket:
\begin{equation*}
    d_{r+1}(X) = \bigsqcup_{y \in \fC_{p-(r+1)}} \<{\fC', \fC_{[p, p - (r+1)]} }_y.
\end{equation*}
\end{thm}
Note that the Toda manifolds are $n-p+r$ dimensional and we get an element in
\begin{equation*}
    \bigoplus_{\norm{y} = p-(r+1)} \Omega^{\fr}_{n-p+r} = \pg{E}{1}{n-1}{p-(r+1)}.
\end{equation*}

\begin{proof}
    First, we construct the chain complex $K = K_{\fC}$ associated to $\fC$ via \cref{constr:FF_to_Ch}. By \cref{lem:ChSur}, there is a spectral sequence with:
    \begin{equation*}
        \pg{E}{1}{n}{p} = \pi_{n - p} K_p = \pi_{n - p} \WedgeSp_{\norm{x} = p} \bbS = \bigoplus_{\norm{x} = p} \Omega^{\fr}_{n - p}.
    \end{equation*}
    
    An extension of $\fC_{[p, p - r]}$ corresponds to a chain complex:
    \begin{equation}\label{eq:long_ch}
        \underbrace{\bbS}_{n+1} \to 0 \to \dots \to 0 \to K_p \to \dots \to K_{p-r}.
    \end{equation}
    where the the manifold $X$ encodes as the higher null homotopies of the map $\bbS \to \dots \to K_p$. Using the cubical model, by \cref{cor:ch_to_morph} we get a morphism of chain complexes
    \begin{equation}\label{eq:long_surv}
        \xymatrix{
            \bbS \ar[r] \ar[d] & 0 \ar[r] \ar[d] & \dotsb \ar[r] \ar[d] & 0 \ar[r] \ar[d] & \dotsb \ar[r] \ar[d] & 0 \ar[d] \\
            0 \ar[r] & \dotsb \ar[r] & \dotsb \ar[r] & K_p \ar[r] & \dotsb \ar[r] & K_{p - r}
        }
    \end{equation}
    see \cref{fig:ext_of_ffcat}. By the universal property of total cofiber this is equivalent to a map of chain complexes:
    \begin{equation*}
        \xymatrix{
            \bbS^{n - p} \ar[d]_X \ar[r] & 0 \ar[r] \ar[d] & \dotsb \ar[r] \ar[d] & 0 \ar[d] \\
            K_p \ar[r] & K_{p - 1} \ar[r] & \dotsb \ar[r] & K_{p - r}
        }
    \end{equation*}
    see \cref{lem:Ch_diff}, in particular iterated the argument about \cref{eq:diff_sus}. The criteria of \cref{lem:ChSur} guarantee that $X$ survive to $E_{r+1}$.
    On the other hand, given such map of chain complexes, we get a chain of the form \cref{eq:long_ch}, and using \cref{prop:FF_are_Ch} we get a lift of this chain complex to a framed flow category which extend $\fC_{[p, p - r]}$.
    
    To determine the differential, recall from \cref{lem:Ch_diff} that the differential $d_r(X)$ is given by the composition of two maps: the extension of $X$, and the map $K_{[p, p - r]}$ viewed as a chain map $K_{[p, p - r + 1]} \to K_{p - r}$. Under the identifications provided by \cref{prop:FF_are_Ch} and \cref{cor:ch_to_morph}, these maps correspond to $\fC'$ and $\fC_{[p, p - r-1]}$, respectively. Therefore, the differential is given by the Toda bracket:
    \begin{equation*}
        d_{r+1}(X) = \langle \fC', \fC_{[p, p - r-1]} \rangle.
    \end{equation*}
    and by decomposing
    \begin{equation*}
        \bigoplus_{\norm{y} = p-(r+1)} \Omega^{\fr}_{n-p+r} = \pg{E}{1}{n-1}{p-(r+1)}.
    \end{equation*}
    we project the differential on each coordinate of the sum
    \begin{equation*}
        d_{r+1}(X) = \bigsqcup_{y \in \fC_{p-(r+1)}} \<{\fC', \fC_{[p, p - r]} }_y.
    \end{equation*}
\end{proof}

\begin{figure}
    \centering
    % https://q.uiver.app/#q=WzAsOCxbMCwwLCIoMCwwLDApPUtfMyJdLFsyLDAsIigwLDAsMSk9S18yIl0sWzIsMiwiKDAsMSwxKT0qIl0sWzMsMywiKDEsMSwxKT1LXzAiXSxbMSwxLCIoMSwwLDApPSoiXSxbMywxLCIoMSwwLDEpPUtfMSJdLFsxLDMsIigxLDEsMCk9KiJdLFswLDIsIigwLDEsMCk9KiJdLFswLDEsIiIsMCx7InN0eWxlIjp7ImJvZHkiOnsibmFtZSI6ImRhc2hlZCJ9fX1dLFsxLDJdLFswLDddLFs3LDIsIiIsMCx7InN0eWxlIjp7ImJvZHkiOnsibmFtZSI6ImRhc2hlZCJ9fX1dLFs0LDZdLFs2LDMsIiIsMix7InN0eWxlIjp7ImJvZHkiOnsibmFtZSI6ImRhc2hlZCJ9fX1dLFs1LDNdLFs0LDUsIiIsMCx7InN0eWxlIjp7ImJvZHkiOnsibmFtZSI6ImRhc2hlZCJ9fX1dLFswLDRdLFsxLDVdLFsyLDNdLFs3LDZdXQ==
    \begin{tikzcd}
        {\bbS^0} & & {*} \\
        & {*} & & {*} \\
        {*} & & {K_1} \\
        & {*} & & {K_0}
        \arrow[from=1-1, to=1-3,                black!30]
        \arrow[from=1-1, to=2-2,                red]
        \arrow[from=1-1, to=3-1,                red!40]
        \arrow[from=1-3, to=2-4,                orange]
        \arrow[from=1-3, to=3-3,                red!40]
        \arrow[from=2-4, to=4-4,                red!40]
        \arrow[from=3-1, to=3-3,                black!30]
        \arrow[from=3-1, to=4-2,                orange]
        \arrow[from=3-3, to=4-4,                red]
        \arrow[from=4-2, to=4-4,                black!30]
        \arrow[from=2-2, to=2-4, crossing over, black!30]
        \arrow[from=2-2, to=4-2, crossing over, red!40]
    \end{tikzcd}
    \caption{A morphism between spine $3$-cubes, which can be seen as a square (light-colored arrows) in the category of spine $2$-cubes (solid arrows). This square is trivial except at the initial and terminal objects (red arrows).}
    \label{fig:ext_of_ffcat}
\end{figure}

\subsection{Gluing manifolds and the differential}

To obtain a more geometric description of the differential in the spectral sequence, we need to define how to glue framed manifolds together. The idea is to construct a closed manifold by gluing together a (punctured) diagram of manifolds.

\begin{defn}
A \textbf{$(n,k)$-boundary} is a diagram $D$ of shape $\op^k \setminus \varnothing$ such that $D(A)$ is a $n-\norm{A}$-dimensional manifold with $\<{k-\norm{A}}$ corners, and the maps are given by embeddings.
\end{defn}

\begin{rem}
Given a spine cube in a category $\cC$, the mapping space is the Permutohedron $P_k$. Under the correspondence between chain complexes and framed flow categories, the Permutohedron corresponds to manifolds with corners of codimension $k - 1$ (the mapping spaces in flow categories). The half Permutohedron $\tfrac{1}{2} P_k$ corresponds to $(k,k)$-boundaries.
\end{rem}

\begin{exmp}
Any manifold with corners of codimension $k$ gives rise to a $(n,k)$-boundary by restricting the $k$-dimensional diagram (as in~\Cref{ex:mfld_diagram}) to $[1]^k \setminus \{ \overline{1} \}$ (and taking $\mathrm{op}$).
\end{exmp}

\begin{lem}\label{lem:TodaManifold}
Let $D$ be a $k$-boundary. Denote by $X = \colim D$ the colimit of $D$, viewed as a diagram valued in topological spaces. Then $X$ has the structure of a smooth $n-1$ dimensional manifold. Moreover, if the $D(A)$ are framed manifolds, then $X$ inherits a framing.
\end{lem}

\begin{proof}
We proceed by induction on $k$. Consider a point $x \in X$ that lies in the image of a manifold $D(A)$ of dimension $n + j$, where $j = \norm{A}$. The point $x$ has a neighborhood of the form $U \subset \bbR^{n + j} \times \partial \Hlf^{k - j + 1}$. If $j > 0$, then by the induction hypothesis, $U$ is homeomorphic to an open subset of $\bbR^{n + j}$. If $j = 0$, then $U$ is homeomorphic to an open subset in $\bbR^{n} \times \partial \Hlf^{k + 1}$.

To smooth $U$, we use stereographic projection. Let:
\begin{equation*}
    H = \left\{ \sum_{i} x_i = 0 \right\} \subset \bbR^{k + 1}
\end{equation*}
and choose a point $p = (\alpha, \alpha, \dotsc, \alpha)$ with $\norm{u} \leq \alpha$ for all $u \in U$. By taking $\alpha$ sufficiently large, the line segment $[p, u]$ intersects $H$, and we define a map:
\begin{equation*}
    s \colon \partial \Hlf^{k + 1} \to H
\end{equation*}
by mapping $u$ to the point of intersection between $[p, u]$ and $H$. The map $s$ is a homeomorphism (showing that $X$ is a topological manifold), and is smooth for $u \neq 0$. Composing the transition functions of the charts of $X$ with $s$ yields smooth transition functions, so $X$ is a smooth manifold.

Since the maps in the diagram $D$ are embeddings, we obtain a map of diagrams:
\begin{equation*}
    D \to BO,
\end{equation*}
where $BO$ is the classifying space for stable vector bundles (treated as a constant diagram). By functoriality of colimits, we have a map:
\begin{equation*}
    c \colon X = \colim D \to \colim BO = BO.
\end{equation*}
The smooth structure on $X$ provides a map $X \to BO$, and we claim that this map is homotopic to $c$. To see this, define homotopies via:
\begin{equation*}
    (t, u) \mapsto t u + (1 - t) s(u),
\end{equation*}
which provides a homotopy between the structure map of $X$ as a smooth manifold and the map $c$.

If the $D(a)$ are framed manifolds, then by functoriality, the map $c$ factors through a point, and thus $X$ inherits a framing.
\end{proof}

\begin{exmp}\label{exmp:boundary_ffcat}\todo[color = lightgray]{redundant duo to \cref{fig:k-boundary}?}
Let $\fC, \sD$ be two framed flow categories with $\fC_{[3,1]} \simeq \sD_{[3,1]}$. Denote by $M$ and $N$ the mapping manifolds of these categories. Under the assumption that:
\begin{equation*}
    M_{j, i} = N_{j, i} \quad \text{for } 3 \geq j, i \geq 1
\end{equation*}
we can define a $3$-boundary $D$ as the following diagram:
\[\begin{tikzcd}
    	& M_{4,3}\times M_{3,2}\times M_{2,1}\times N_{1,0} \\
    	{M_{4,3}\times M_{3,2}\times N_{2,0}} & {M_{4,3}\times M_{3,1}\times N_{1,0}} & {M_{4,2}\times M_{2,1}\times N_{1,0}} \\
    	{M_{4,3}\times N_{3,0}} & {M_{4,2}\times N_{2,0}} & {M_{4,1}\times N_{1,0}} \\
    	\arrow[from=1-2, to=2-1]
    	\arrow[from=1-2, to=2-2]
    	\arrow[from=1-2, to=2-3]
    	\arrow[from=2-1, to=3-1]
    	\arrow[from=2-2, to=3-1]
    	\arrow[from=2-2, to=3-3]
            \arrow[from=2-3, to=3-3]
            \arrow[from=2-1, to=3-2, crossing over]
    	\arrow[from=2-3, to=3-2, crossing over]
    \end{tikzcd}\]
The framed manifold $X = \colim D$ is the obstruction to lifting $D$ to a framed flow category.
\end{exmp}

We now generalize \cref{exmp:boundary_ffcat}.

\begin{defn}\label{def:GeoToda}
Let $\fC,\sD$ be framed flow categories and $n \in \bbZ$ such that $\fC_{k}$ is trivial for  $n+1 >  k > p-r$, and $\fC_{[p,p-r]} \simeq \sD_{[p,p-r]}$.  Denote by $M$ and $N$ the mapping manifolds of $\fC$ and $\sD$. For any $y \in \sD_{p-r-1}$ we define the $r+1$-boundary diagram
\begin{equation*}
    D_y \colon  P(\{p, \dots, p-r\})\setminus \varnothing \to \Top
\end{equation*}
in the following way:
\begin{equation*}
    D(a_1,\dots, a_k) \coloneqq M_{p+1,a_1} \times \dotsm \times M_{a_{k-1}, a_k} \times N_{a_k, y}.
\end{equation*}
\end{defn}

\begin{lem}\label{lem:Toda_FFCat_is_Colim_Manifold}
Let $\fC, \sD$ as above, Then:
\begin{equation*}
    \<{\fC, \sD}_y = \colim D_y
\end{equation*}
as elements in $\Omega^{\operatorname{fr}}_{r}$.
\end{lem}

\begin{proof}
Denote by $K_{\fC} \in \Fun(\J_{[n + 1, p-r]}, \Sp)$ and $K_{\sD} \in \Fun(\J_{[p,p-r-1]}, \Sp)$ the chain complexes corresponding to $\fC_{[n + 1, p-r]}$ and $\sD_{[p,p-r]} \sqcup \{y\}$, respectively.

By the definition of the Toda bracket (of flow categories) as the composition of chain maps, and using~\Cref{lem:diff_is_comp}, we find that $\langle \fC, \sD \rangle_y$ is given by the map:
\begin{equation}\label{eq:OneOb}
    \bigcup_{n + 1 \geq k \geq 1} (\Sinf \Hlf^{\{1, \dotsc, k - 1\}})_+ \smashSp (\Sinf \Hlf^{\{k + 1, \dotsc, n + 1\}})_+ \to \bbS,
\end{equation}
where the equivalence $\Cu_{\mathbb{Z}} \simeq \J$ allows us to move from the mapping spaces in $\Cu_{\mathbb{Z}}$ (the Permutohedron) to mapping spaces in $\J$.

We can think of this union as the colimit over subsets $A \subset J_{n + 1, 0}$ and we denote:
\begin{equation*}
    \partial (\Hlf^{n + 2})_+ \coloneqq (\Hlf^{n + 2})_+ \setminus \{ \overline{1} \},
\end{equation*}
where we treat $\Hlf^{n + 2}$ as an $(n + 2)$-diagram. Then, we can rewrite the Toda bracket as:
\begin{equation*}
    \alpha\colon\colim \partial (\Sinf \Hlf^{n + 2})_+ \to \bbS.
\end{equation*}

The restriction of $\alpha$ to the $k$ boundary component is denoted:
\begin{equation*}
    \alpha_k \colon \Sigma^\infty_+ \Hlf^{\{1, \dotsc, k - 1\}} \smashSp \Sigma^\infty_+ \Hlf^{\{k + 1, \dotsc, n + 1\}} \to \bbS.
\end{equation*}
Using the definitions of $K_{\fC}$ and $K_{\sD}$, we have that $\alpha_k$ are the collapse maps:
\begin{equation*}
    \alpha_k = c_{M_{n + 1, k + 1}} \smashSp c_{N_{k, 0}} = c_{M_{n + 1, k + 1} \times N_{k, 0}}.
\end{equation*}

The maps in the diagram $\partial (\Hlf^{n + 2})_+ \to \bbS$ are given by the $\alpha_k$ and higher compositions involving more than two manifolds.

Denote:
\begin{equation*}
    W = \colim D.
\end{equation*}
The framing on $W$ is given by taking the colimit of the framings of the $D(A)$'s. Combining this with the contractibility of the embedding spaces, we find that the Pontryagin--Thom collapse map of $W$ is the colimit:
\begin{equation*}
    c_W \simeq \colim_A c_{D(A)}.
\end{equation*}
We have:
\begin{equation*}
    \alpha_k = c_{D(\{ k \})}.
\end{equation*}
Thus, we conclude that the Toda bracket is given by:
\begin{equation*}
    \colim \partial (\Hlf^{n + 2})_+ \xrightarrow{c_{D(A)}} \bbS,
\end{equation*}
which is precisely $c_W$.

Applying the Pontryagin--Thom isomorphism completes the proof.
\end{proof}


Combining the geometric description of the pages and the differentials, we obtain the following result.

\begin{cor}\label{cor:SSofFFCat}
Let $\fC$ be a framed flow category. Then there exists a spectral sequence converging to the realization
\begin{equation*}
    \rnp{E} \implies \pi_n \norm{\fC},
\end{equation*}
with:
\begin{equation*}
    \pg{E}{1}{n}{p} = \bigoplus_{\norm{y} = p} \Omega^{\fr}_{n - p}.
\end{equation*}
An element $X \in \pg{E}{1}{n}{p}$ survives to the $(r+1)$-th page if and only if there is an $(X,r,n,p)$ extension of $\fC$.
Let $\fC'$ be such an extension, then the differential is represented in $\pg{E}{1}{n-1}{p-(r+1)}$ by the \emph{Toda manifold}
\begin{equation*}
    d_{r+1}(X) = \bigsqcup_{y \in \fC_{p-r-1}} \colim D_y,
\end{equation*}
viewed as an element in $\Omega^{\fr}_{n - p + r}$.
\end{cor}


\subsection{Application to Morse homotopy theory}

The work of Cohen, Jones, and Segal \cite{CJS95}, together with the framework of \cite{Qin18}, connects Morse theory, flow categories, and the structure of chain complexes. The results of this section can be seen as an application of these ideas.

\begin{defn}\label{def:FlowMan}
Let $M$ be a closed smooth manifold (without corners) equipped with a Smale-Morse function $f \colon M \to \bbR$. For critical points $a, b$ of $f$, denote by $\phi$ the gradient flow of $-f$, and define to be the moduli space of flow lines from $a$ to $b$ modulo reparametrization:
\begin{equation*}
    M(a, b) = \left\{ x \in M \mid \lim_{t \to -\infty} \phi_t(x) = a, \ \lim_{t \to \infty} \phi_t(x) = b \right\} / \bbR,
\end{equation*}
,where the $\bbR$ action is given by $r.x = \phi_r(x)$.
\end{defn}

\begin{thm}[Cohen--Jones--Segal \cite{CJS95}, Qin \cite{Qin18}]\label{thm:Morse_FFcat}
There exists a framed flow category $\fC_f$, where the objects are the critical points of $f$, with grading given by the Morse index $\norm{x}$ of $x$, and the morphism spaces are the compactified moduli spaces:
\begin{equation*}
    \Hom_{\fC_f}(x, y) = \overline{M}(x, y),
\end{equation*}
where $\overline{M}(x, y)$ is the compactification of $M(x, y)$. The composition is given by concatenation of flow lines, and the framings come from the stable and unstable manifolds of the critical points.
\end{thm}

\begin{thm}[Cohen--Jones--Segal \cite{CJS95,Cohen19}]\label{thm:realization}
The realization (total cofiber) of the chain complex $K_f$ associated to the flow category $\fC_f$ is the suspension spectrum of $M$:
\begin{equation*}
    \norm{\fC_f} \simeq \Sigma^\infty_+ M.
\end{equation*}
\end{thm}

Using the results of the previous sections, we can construct a spectral sequence for the Morse function flow category.

\begin{thm}\label{thm:SSofMorseFunction}
Let $M$ be a smooth Riemannian manifold, and let $f \colon M \to \bbR$ be a Morse--Smale function on $M$. For a critical point $x$ of $f$, denote by $\norm{x}$ the Morse index of $x$. Then there exists a spectral sequence with:
\begin{equation*}
    E^1_{p, n} = \WedgeSp_{\norm{x} = p} \Omega^{\operatorname{fr}}_{n - p},
\end{equation*}
converging to $\Sigma^\infty_+ M$ (the stable homotopy type of $M$).

Denote by $\fC_f$ the flow category of $(M, f)$. Given an element $X \in E^1_{p, n}$, we have that $X$ survives to the $(r+1)$-th page if and only if there is an $(X,r,n,p)$ extension of $\fC_f$. Let $\fC'$ be such an extension. Then the differential is represented in $\pg{E}{1}{n-1}{p - (r+1)}$ by the Toda bracket
\begin{equation*}
    d_{r+1}(X) = \bigsqcup_{y \in \fC_{p-(r+1)}} \langle \fC', \fC_{[p, p - r]} \rangle_y.
\end{equation*}
\end{thm}

\begin{proof}
This follows directly from \cref{thm:SSofFFCat} and \cref{thm:realization}, applying them to the flow category $\fC_f$ associated to the Morse function $f$.
\end{proof}